% 
\section{Related Works}
\label{sec:rel}

% 
\subsection{Deep Learning-based Visual SLAM}
\label{sec:rel1}

Following traditional modular visual SLAM systems \cite{murTRO2015,klein2007parallel}, research integrating deep learning into the SLAM pipeline has become mainstream \cite{favorskaya2023deep,teed2021droid,zhu2022nice}. A prominent example of this trend is \text{DROID-SLAM} \cite{teed2021droid}, which set a new standard by achieving high accuracy through a fully differentiable architecture. More recently, a subsequent wave of \text{feed-forward dense SLAM} systems has emerged, which leverage pre-trained two-view reconstruction foundation models. \text{MASt3R-SLAM} \cite{murai2025mast3r} and \text{VGGT-SLAM} \cite{maggio2025vggt}, which attempt to address the ambiguity of the scale, are notable examples of this approach. The aforementioned works aim to maintain the robustness of modern learning-based \ac{SLAM} systems in real-world environments. However, their robustness to scale drift, particularly in large-scale and complex indoor environments, remains largely unexplored.

% 
\subsection{Benchmark Datasets for Visual SLAM}
\label{sec:rel2}
The progress and evaluation of visual SLAM algorithms have been critically dependent on the role of standard benchmark datasets that provide precise ground truth. The rigorous performance assessment in SLAM research is primarily based on three key datasets. \text{TUM-RGBD} dataset \cite{sturm2012benchmark}  established the standard for evaluating \ac{ATE}, while the \text{EuRoC MAV} dataset \cite{burri2016euroc} has served as a crucial benchmark for the robustness of visual-inertial systems. Lastly, \text{7-Scenes} \cite{glocker2013real} is specialized for evaluating re-localization performance. These are collectively regarded as the de facto standards for performance validation in the visual SLAM community. Despite their utility, standard benchmarks are mostly room-scale and do not support systematic evaluation of long-term \textbf{intra-session} scale drift or \textbf{inter-session} scale consistency. ARKitScenes \cite{dehghan2021arkitscenes}—though a comparatively recent dataset—also falls short, as it does not encompass complex indoor environments with long trajectories. This lack of a scale-focused benchmark has limited systematic analysis of scale failure modes; this is a gap we address in this work.

% % 
% \subsection{Evaluation Metrics for Visual \ac{SLAM}}
% \label{sec:rel3}
% A commonly reported trajectory metric is the \ac{ATE} \cite{Zhang18iros}, computed after Sim(3) alignment of the estimated poses to ground truth; it summarizes global trajectory agreement as an \ac{RMSE} over pose errors. While invaluable, \ac{ATE} only assesses the accuracy of the localization component (L in \ac{SLAM}). It provides no direct information about the quality of the mapping component (M), which is the primary output of dense \ac{SLAM} systems \cite{li2023dense}.
% To evaluate 3D structures, the computer vision community commonly employs metrics like Chamfer distance, which measures the average distance between two point clouds. Such metrics provide a direct assessment of geometric similarity and have been used for tasks like point cloud completion. However, their application to evaluating monocular dense visual SLAM systems has been limited by a scarcity of high-fidelity 3D ground truth maps. While several recent datasets \cite{wei2025fusionportablev2,chen2024heterogeneous} provide 3D ground-truth maps for evaluating reconstruction quality, most of them focus on outdoor environments. Indoor datasets such as 7-Scenes \cite{glocker2013real} or ScanNet++ \cite{yeshwanth2023scannet++} exist, but they are confined to room-scale scenarios and thus remain insufficient for revealing scale inconsistency in large and complex indoor environments.
% % Specifically, some 3D datasets exist, they are often confined to room-scale scenarios.

% Our work is the first to bridge this critical gap by providing a large-scale benchmark with 3D ground truth specifically designed for this purpose, and by incorporating Chamfer distance into a comprehensive SLAM evaluation protocol.


\subsection{Evaluation Metrics for Visual \ac{SLAM}}
\label{sec:rel3}
A commonly reported trajectory metric is the \ac{ATE} \cite{Zhang18iros}, computed after Sim(3) alignment of the estimated poses to ground truth; it summarizes global trajectory agreement as a \ac{RMSE} over pose errors. While invaluable, \ac{ATE} only assesses the accuracy of the localization component (L in \ac{SLAM}). It provides no direct information about the quality of the mapping component (M), which is the primary output of dense \ac{SLAM} systems \cite{li2023dense}. To directly assess 3D structure quality, we adopt the Chamfer distance \cite{murai2025mast3r} (and Drop Rate will be defined in \secref{sec:exp-setup}) between the reconstructed map and a high-fidelity reference point cloud. These map-oriented metrics provide geometry-focused signals that complement trajectory-only evaluation and have been used in prior dense visual SLAM work. Although many \ac{SLAM} datasets exist (e.g., \cite{wei2025fusionportablev2,chen2024heterogeneous}), the benchmarks that deep, dense monocular visual SLAM systems have been primarily evaluated on are those highlighted in \secref{sec:rel2} (e.g., \cite{sturm2012benchmark,burri2016euroc,glocker2013real}). These room-scale or structurally simple indoor datasets do not capture challenging large-volume environments (e.g., multi-floor lobbies with tens of meters of ceiling height), where long-range intra- and inter-session scale inconsistency becomes apparent. Because dense visual SLAM yields dense maps as well as poses, map-oriented evaluation is crucial: ATE can look correct while the map suffers scale distortions. With Chamfer- and Drop Rate–based evaluation on ScaleMaster, we provide a protocol that exposes such failures under conditions existing benchmark datasets cannot reveal.

% To evaluate 3D structures, the computer vision community commonly employs metrics such as Chamfer distance, which measures the average distance between two point clouds. These metrics provide a direct assessment of geometric similarity and have been widely used in tasks like point cloud completion and 3D reconstruction. However, their application to monocular dense visual SLAM has been limited by the scarcity of high-fidelity 3D ground truth maps. While several recent datasets \cite{wei2025fusionportablev2,chen2024heterogeneous} provide 3D LiDAR point clouds, they are mostly designed for sensor fusion in outdoor environments rather than as ground truth maps for dense visual SLAM evaluation. Indoor datasets \cite{guo2024lidar,yeshwanth2023scannet++} exist, but they are confined to room-scale scenarios and are thus insufficient for exposing scale inconsistency in large and complex indoor environments.

% We adopt Chamfer distance as our primary map-to-map evaluation metric not only because it provides a direct and geometry-focused measure of fidelity, but also because it has been used in prior dense visual SLAM work such as MASt3R-SLAM, enabling fair and direct comparison with existing results. Our work is the first to bridge this critical gap by providing a large-scale indoor benchmark with 3D ground truth specifically designed for this purpose, and by incorporating Chamfer distance into a comprehensive SLAM evaluation protocol.
